%% Generated from paper/draft/06_prescription.md by md2tex.py.
%% Mechanical first pass -- polish by hand, then stop regenerating.

\section{Compile-time configuration}
\label{sec:06}

\subsection{Why this is the remaining lever}

Section~\ref{sec:05} closed the runtime path by a branch defined before the computation was run: if generation-length or return-time knowledge had recovered at least half the headroom, a prediction-based runtime policy would have remained a candidate. It recovered at most a small fraction, so \textbf{compile-time configuration is what is left} --- coordination decided once, in advance, for everybody. This verdict is about this substrate under a preregistered branch; a server-side scheduler remains an open path that this work did not take.
% CLAIMS 3.1
% CLAIMS 2.6

\subsection{Choosing a configuration without touching the device}

Candidate configurations are scored by the uncalibrated model on \textbf{exploration} seeds and confirmed on \textbf{evaluation} seeds. The only workload facts entering the choice are the gap distribution and the concurrency ceiling; no device measurement enters the selection loop. A sensitivity analysis says the choice is insensitive to the gap distribution and responds only to the concurrency ceiling --- \textbf{that is the one statistic worth re-measuring before reconfiguring.}
% CLAIMS 3.2
% CLAIMS 3.9

The cost of acting on the choice is a recompile, and it is small: compile time is linear in the number of compiled graphs. The model was fitted to the first two observations alone; at the fifth, sixth and seventh points its predictions held to within +3.0 \% on compile time and +1.6 \% on artifact size. On this model and device a six-rung artifact takes about eight minutes.
% CLAIMS 3.3

\begin{figure}[t]
  \centering
  \includegraphics[width=\columnwidth]{fig7_compile_cost}
  \caption{A recompile is cheap enough to treat as a configuration choice, and its cost is predictable. Wall-clock against the number of compiled graphs; the dashed line is a model fitted to the two square points alone, and the circles are the later observations that tested it. Qwen3-4B at a sequence length of 8{,}192 on four devices.}
  \label{fig:fig7}
\end{figure}


\subsection{Confirmation on device}

Predictions were committed before measurement, together with the judgement rule and the channel definition. Against a baseline of \texttt{batch\_size = 8} on grid \texttt{(1,2,4,8)}, a tuned configuration (\texttt{batch\_size = 16}, grid \texttt{(1,4,6,8,10,16)}) recovered \textbf{+9.72 \% / +10.07 \%} of device time on the two channels at N=8 concurrent sessions, and \textbf{+2.11 \% / +2.74 \%} at N=6 on a fresh seed under a corrected channel-agreement rule.
% CLAIMS 3.4
% CLAIMS 3.5

A device-side ablation separates the two axes. Changing only \texttt{batch\_size} yields +8.25 \%; adding grid alignment on top adds a further 1.3--1.5 percentage points. The two arms' \textbf{prefill ratios are identical in measurement} (0.678 / 0.678), which is what makes the ablation interpretable: with the same pool size the cache survives identically, so the arms can differ only in the decode term. The causal chain closes: \texttt{batch\_size} sets the KV block count, that sets cache survival, and survival sets prefill recomputation.
% CLAIMS 3.6, 3.7
% CLAIMS 3.8

\begin{figure}[t]
  \centering
  \includegraphics[width=\columnwidth]{fig8_final_result}
  \caption{A configuration chosen without any device measurement recovers device time, and the source of the gain is conditional rather than fixed. Bars are channel A$'$, hollow circles channel B, crosses the predictions committed before measurement; the lighter portion of each right-hand bar is what grid alignment adds on top of pool size. At N=8 the pool dominates; at N=6 on a fresh seed the baseline already reuses 17 of 18, so almost nothing is left for the pool to buy.}
  \label{fig:fig8}
\end{figure}


\subsection{The gain is proportional to what is left to recover}

The ablation's ordering held on a fresh seed but the \textbf{magnitudes reversed}: where the baseline already reused 17 of 18 resumptions, changing only \texttt{batch\_size} bought +0.59 \% --- inside the equivalence band --- and nearly all of the remaining gain came from grid alignment instead. The reading of Section~\ref{sec:06} is therefore not "\texttt{batch\_size} dominates" but \textbf{"the dominant factor is recoverable cache loss, and \texttt{batch\_size} is the lever that buys it."} The model predicted this reversal before measurement.
% CLAIMS 3.10

\subsection{Where the lever stops}

% 편집 주석: [SWEEP-PENDING] 자리. TASK40으로 채움.

If a larger pool recovers cache, the obvious rebuttal is to make it much larger. We tested it directly: \texttt{batch\_size} $\in$ \{8, 16, 24, 32\}, with the grid held at \texttt{(1,4,6,8,10,B)} so that \textbf{only the top rung changes} between adjacent configurations, over three concurrency levels and three blocks.

The gain \textbf{saturates at 16}. Judged by the bootstrap CI of the median per-cell ratio, both controlled steps are equivalence: 16 $\rightarrow$ 24 gives 0.9999 (CI [0.9808, 1.0000]) and 24 $\rightarrow$ 32 gives 1.0002 (CI [0.9999, 1.0015], width 0.0016) --- at N $\leq$ 10 under this workload.
% CLAIMS 3.11

Two mechanisms make the plateau, and the data shows both. \textbf{Survival saturates}: at \texttt{batch\_size = 16} reuse already reaches 18/18 and 24/24 at the two lower concurrencies, so a larger pool has no cache left to save. \textbf{And the widest rung is never selected}: the mapping table places rung 16, 24 and 32 at a concurrency of eleven, which this grid never reaches, and the step trace confirms a 0.0 \% share for the top rung in every cell. A rung that is never selected can neither save nor cost anything --- \textbf{without this second mechanism the curve would have reversed}, since the extrapolated step cost of the widest rung is roughly twice that of rung 8.
% CLAIMS 3.10
% CLAIMS 3.12

The plateau leaks in exactly one place, and it leaks where the law of Section~\ref{sec:06} says it should. At N=10 the baseline pool reached only 28/30 survival, and there --- and only there --- 16 $\rightarrow$ 24 still bought about 2 \%.
% CLAIMS 3.10

The practical consequence is the point of the section. Device memory on this instance is \texttt{2.1 + 0.28125$\cdot$B} GiB per device (a three-point fit), so the 15.7 GiB ceiling is reached near \texttt{B = 46} --- \textbf{an extrapolation, not a measurement}. The gain ends at one third of that. \textbf{The optimal \texttt{batch\_size} is set by the workload distribution, not by what the hardware can hold.}
% CLAIMS 3.13

One further result belongs here because of what it tests. The uncalibrated model predicted this plateau \textbf{before measurement}, and at the two lower concurrencies its six cells were accurate to 0.0003--0.0026 --- the tightest agreement in this project. Predicting the \emph{absence} of an effect is a different kind of evidence from predicting its size: there is little room for a null prediction to be right by accident.
% CLAIMS 3.14

\begin{figure}[t]
  \centering
  \includegraphics[width=\columnwidth]{fig9_batch_saturation}
  \caption{The gain from a larger KV pool ends where the survival rate saturates, at one third of what the device can physically hold. Upper panel is device time against batch\_size with the grid held at $(1,4,6,8,10,B)$ so that only the top rung changes; lower panel is layer-2 survival on the same axis. The vertical dashed line is the KV ceiling extrapolated from a three-point fit of device memory, not a measurement. The circled point is the one cell whose survival had not yet saturated, and the only one where a larger pool still bought anything. Concurrency of ten and below.}
  \label{fig:fig9}
\end{figure}


\subsection{The prescription}

Putting Section~\ref{sec:06}--Section~\ref{sec:06} together: \textbf{measure the workload's concurrency ceiling, choose the smallest pool at which the cache survives at that ceiling, and stop.} Grid alignment is a real but secondary gain on top. Nothing in this procedure requires predicting an individual tool call --- only its distribution.
